\documentclass[11pt,a4paper]{article}
\usepackage[margin=28mm]{geometry}
\usepackage[T1]{fontenc}
\usepackage{parskip}
\usepackage[hidelinks]{hyperref}
\begin{document}
\begin{flushright}
26 September 2026
\end{flushright}

Professor K. Ravi-Chandar\\
Editor-in-Chief\\
\textit{International Journal of Fracture}

Dear Professor Ravi-Chandar,

On behalf of my co-authors, I am pleased to submit the manuscript
\textbf{\textquotedblleft Analytical Model for Branching Initiation from the Corner of a
V-Shaped Interfacial Shear Crack\textquotedblright} by Anatoly Kaminsky, Mikhailo Selivanov,
Mykhailo Dudyk, and Tetiana Polishchuk for consideration as a Research article
in the \textit{International Journal of Fracture}.

The manuscript develops an analytical fracture-mechanics model for branching
initiation from the corner of a prescribed stationary V-shaped interfacial
shear crack located along a kinked bimaterial interface. The singular corner
field of the stationary crack is used as the outer field for a small tensile
process zone in one of the adjoining elastic materials. By combining Mellin
transformation, asymptotic matching, Wiener--Hopf factorization, and stress
regularity at the process-zone tip, we obtain closed-form relations for the
equilibrium zone length, opening, work of separation, and a zone-related
energy parameter. The analysis also establishes the physically admissible
material/loading branches, quantifies the effects of interface angle and
elastic contrast, and distinguishes corner branching from competing
interfacial propagation at the crack-branch tips.

We believe the manuscript fits the scope of the \textit{International Journal
of Fracture} because it is an analytical contribution aimed at improving the
mechanistic understanding of fracture initiation in a bimaterial system with
a non-classical crack geometry.

The present work is related to, but distinct from, several earlier studies.
Kaminsky, Dudyk, and Chornoivan (\textit{International Journal of Fracture}, vol. 250, Article 22, 2026) considered the influence of T-stress on process-zone modelling
for kinking from the tip of a conventional interfacial crack. The present
manuscript addresses a different local fracture configuration: branching from
the corner of a stationary V-shaped interfacial shear crack, for which the
governing outer field is a corner singularity rather than the usual
interfacial crack-tip field. Nazarenko and Kipnis
(\textit{International Applied Mechanics}, vol. 58, no. 5, pp. 497--509, 2022) analysed the limiting
equilibrium and interfacial propagation of the V-shaped shear-crack branches;
here that stationary-crack solution provides the outer asymptotic field for a
new local process-zone problem at the corner.

The manuscript is original, has not been published previously, and is not
under consideration elsewhere. All authors have approved the submitted
version and agree with its submission to the \textit{International Journal
of Fracture}.

Thank you for considering our manuscript.

Sincerely,

\textbf{Mikhailo Selivanov}\\
Corresponding author\\
S. P. Timoshenko Institute of Mechanics\\
National Academy of Sciences of Ukraine\\
Kyiv, Ukraine\\
\href{mailto:mfs@ukr.net}{mfs@ukr.net}
\end{document}
